%!TEX root = ../template.tex
%%%%%%%%%%%%%%%%%%%%%%%%%%%%%%%%%%%%%%%%%%%%%%%%%%%%%%%%%%%%%%%%%%%%
%% chapter2.tex
%% NOVA thesis document file
%%
%% Chapter with lots of dummy text
%%%%%%%%%%%%%%%%%%%%%%%%%%%%%%%%%%%%%%%%%%%%%%%%%%%%%%%%%%%%%%%%%%%%

\typeout{NT FILE chapter2.tex}%

\chapter{Problem Formulation}
\label{ch:problem_formulation}

\paragraph*{Notation} 
\par The set of real numbers is denoted by $\mathbb{R}$. Given a vector $\vec{v} \in \mathbb{R}^n$, its transpose is denoted by $\vec{v}'$, and likewise for a matrix $\mathbf{M}$. The operation diag($\vec{v}$) creates a matrix with diagonal elements v and off-diagonal elements 0. The inner product between two matrices or vectors is denoted by $\vec{v} \cdot \mathbf{M}$. The cross product operation between a vector and a matrix, denoted by $\vec{v} \times \mathbf{M}$, is equivalent to the inner product between the skew asymmetric matrix of that vector and that matrix, which can be denoted by $[\vec{v}] \cdot \mathbf{M}$. The operation vec($\mathbf{M}$) returns the vectorized matrix.
\par When expressing the coordinates of a point or vector on an object, it may be expressed in the inertial reference frame, denoted by $\mathcal{F}_I = (\vec{e}_I, \vec{e}_J, \vec{e}_K)$, which refers to the coordinate system whose origin is fixed on the ground, or the body reference frame, denoted by $\mathcal{F}_B = (\vec{e}_x, \vec{e}_y, \vec{e}_z)$, which has it's origin on the center mass of the object. When a vector has an uppercase B in overscript ($\vec{v}^B$), it's expressed in the body reference frame. In the absence of an overscript notation, assume it's expressed in the inertial reference frame. \\

\section{Modeling}
Designing the quad-rotor's flight controller will require both vehicle's mathematical models that capture their flight dynamics as accurately as possible. 

\subsection{Missile}
\label{subsec:missile}
The missile's model is built on the assumptions that it's in the midflight guidance phase, mass steadily decreases as fuel is spent, method of propulsion is a two-pulse rocket motor that provides constant downward acceleration (in $\mathcal{F}_B$), the attitude actuators are a ring of pulse jets near the top, it has GNC capability that allows it to discern its position relative to obstacles and the target and path plan, and it is subject to aerodynamic drag. \\

The missile's model is represented by the state vector $x_R = [\vec{p}_{R}, \vec{v}_{R}, \vec{\omega}^B, \text{vec}(\mathbf{A})]$, considered to be a point-mass at the missile's center of mass, in which
\begin{description}
    \item[$\vec{p}_{R} \in \mathbb{R}^3$] is the position of the center of mass, expressed in $\mathcal{F}_I$,
    \item[$\vec{v}_{R} \in \mathbb{R}^3$] is the linear speed of the center of mass, expressed in $\mathcal{F}_I$,
    \item[$\vec{\omega}^B \in \mathbb{R}^3$] is the angular velocity of the rocket, expressed in $\mathcal{F}_B$, and
    \item[$\mathbf{A} \in \mathbb{R}^{3 \times 3}$] is Euler's rotational matrix of the rocket, expressed in $\mathcal{F}_I$. \\
\end{description}

Additional variables are defined as
\begin{description}
    \item[$\vec{a}^B \in \mathbb{R}^3$] is the thruster acceleration vector acting on the rocket, expressed in $\mathcal{F}_B$,
    \item[$\vec{g} \in \mathbb{R}^3$] is the gravitational acceleration vector, expressed in $\mathcal{F}_I$,
    \item[$\mathbf{J} \in \mathbb{R}^{3 \times 3}$] is the rocket's inertia matrix, and
    \item[$\vec{\tau}^B \in \mathbb{R}^3$] is the torque vector on the rocket body, expressed in $\mathcal{F}_B$.\\
\end{description}

Furthermore, $\mathbf{J}$, $\vec{\tau}^B$ and $\vec{a}^B$ are given by the expressions
\begin{equation*}
    \vec{\tau}^B = \begin{bmatrix}
        0 \\
        0 \\
        -L \\
    \end{bmatrix} 
    \times (m \cdot \vec{a}^B),
\end{equation*}

\begin{equation*}
    \mathbf{J} = \text{diag}(\frac{1}{4}m \cdot R^2 + \frac{1}{12} m (2L)^2, \quad \frac{1}{4}m \cdot R^2 + \frac{1}{12} m (2L)^2, \quad \frac{1}{2}m \cdot R^2),
\end{equation*}

\begin{equation*}
    \vec{a}^B = \begin{bmatrix}
        a_x^B \\
        a_y^B \\
        a_z^B \\
    \end{bmatrix}.
\end{equation*}

The lateral acceleration input vector $u_R \in \mathbb{R}^2$ is defined as
\begin{equation*}
    u_R = \begin{bmatrix}
        a_x^B \\
        a_y^B \\
    \end{bmatrix}
\end{equation*}

Figure~\ref{fig:missile_diag} shows the missile's state variables, apart from $\mathbf{A}$, in their respective referential. Even though, mathematically speaking, each vector actually originates from the origin of their referential, here they have been represented as originating from their "point of application" for a more intuitive viewing. \\

\begin{figure}[htbp]
  \centering
    \subbottom[Missile's state variables in $\mathcal{F}_I$]{%
    \label{fig:missile_diagIJK}
    \includegraphics[width=0.45\textwidth]{ch2_missile_diagIJK}
    }
    \subbottom[Missile's state variables in $\mathcal{F}_B$]{%
    \label{fig:missile_diagIJK}
    \includegraphics[width=0.45\textwidth]{ch2_missile_diagXYZ}
    }
  \caption{Missile's state vector diagram}
  \label{fig:missile_diag}
\end{figure}

The variable's dynamics are thus defined by the equations
\begin{align*}
    \dot{\vec{p}}_{R} &= \vec{v}_{R} = f_p(x_R), \\
    \dot{\vec{v}}_{R} &= \mathbf{A} \cdot \vec{a}^B + \vec{g}  = f_v(x_R, u_R), \\
    \dot{\vec{\omega}}^B &= \mathbf{J}^{-1} \cdot (\vec{\tau}^B - \vec{\omega}^B \times \mathbf{J} \cdot \vec{\omega}^B) = f_{\omega}(x_R, u_R), \\
    \text{vec}(\dot{\mathbf{A}}) &= \text{vec}(-\vec{\omega}^B \times \mathbf{A}) = f_A(x_R). \\
\end{align*}

The missile's differential dynamics are defined as $\dot{x}_{R} = f_R(x_R, u_R)$, and 

\begin{equation*}
    f_R(x_R, u_R) = \begin{bmatrix}
        f_p(x_R) \\
        f_v(x_R, u_R) \\
        f_{\omega}(x_R, u_R) \\
        f_A(x_R) \\
    \end{bmatrix}.
\end{equation*}

\begin{comment}
\begin{align*}
    \dot{x}_{R} = f_R(x_R, u_R) = \begin{bmatrix}
        f_p(x_R, u_R) \\
        f_v(x_R, u_R) \\
        f_{\omega}(x_R, u_R) \\
        f_A(x_R, u_R) \\
    \end{bmatrix} = 
    \begin{bmatrix}
        \vec{v}_{R} \\
        \mathbf{A} \cdot \vec{a}^B + \vec{g}\\
        \mathbf{J}^{-1} \cdot (\vec{\tau}^B - \vec{\omega}^B \times \mathbf{J} \cdot \vec{\omega}^B) \\
        \text{vec}(-\vec{\omega}^B \times \mathbf{A}) \\
    \end{bmatrix} = 
    \begin{bmatrix}
        \dot{\vec{p}}_{R} \\
        \dot{\vec{v}}_{R} \\
        \dot{\vec{\omega}}^B \\
        \text{vec}(\dot{\mathbf{A}}) \\
    \end{bmatrix}
\end{align*}

\end{comment}

\subsection{Quad-Rotor Drone}
The quad-rotor drone's \gls{6DoF} model is built on the assumptions that its mass remains constant, the method of propulsion are four vertical rotors, attitude control is provided by a torque vector and altitude control by a scalar thrust input, it has GNC capability that allows it to discern its position relative to obstacles and the target and path plan, and it is subject to aerodynamic drag. The model is based on the Euler-Angle-Free model described in ~\cite{AssloujE}. \\

The drone's model is represented by the state vector $x_D = [\vec{p}_{D}, \vec{v}_{D}, \vec{\omega}^B, \text{vec}(\mathbf{A})]$, considered to be a point-mass at the drone's center of mass. Furthermore,\\

\begin{description}
    \item[$\vec{p}_{D} \in \mathbb{R}^3$] is the position of the center of mass, expressed in $\mathcal{F}_I$,
    \item[$\vec{v}_{D} \in \mathbb{R}^3$] is the linear speed of the center of mass, expressed in $\mathcal{F}_I$,
    \item[$\vec{\omega}^B \in \mathbb{R}^3$] is the angular velocity of the drone, expressed in $\mathcal{F}_B$, and
    \item[$\mathbf{A} \in \mathbb{R}^{3 \times 3}$] is Euler's rotational matrix of the drone, expressed in $\mathcal{F}_I$. \\
\end{description}

Additional variables are defined as
\begin{description}
    \item[$\vec{g} \in \mathbb{R}^3$] is the gravitational acceleration vector, expressed in $\mathcal{F}_I$,
    \item[$\eta_v, \eta_{\omega} \in \mathbb{R}$] are the drag coefficient for the linear and angular velocities, respectively,
    \item[$F \in \mathbb{R}$] is the absolute value of the total thrust given by the rotors, 
    \item[$\mathbf{J} \in \mathbb{R}^{3 \times 3}$] is the drone's inertia matrix, and
    \item[$\vec{\tau} \in \mathbb{R}^3$] is the torque vector on the drone body, expressed in $\mathcal{F}_I$.\\
\end{description}

Furthermore, $\mathbf{J}$, is given by the expression

\begin{equation*}
    \mathbf{J} = \text{diag}(-\frac{1}{3}(2R^2m(\frac{R^3}{8} - (L + \frac{R}{2})^3)), \quad
    -\frac{1}{3}(2R^2m(\frac{R^3}{8} - (L + \frac{R}{2})^3)), \quad
    \frac{1}{12}(R^4m(2L + R)))
\end{equation*}

The thrust and torque inputs are in the vector $u_D \in \mathbb{R}^4$ which is defined as
\begin{equation*}
    u_R = \begin{bmatrix}
        F \\
        \tau \\
    \end{bmatrix}
\end{equation*}

The variable's dynamics are thus defined by the equations
\begin{align*}
    \dot{\vec{p}}_{D} &= \vec{v}_{D} = f_p(x_D), \\
    \dot{\vec{v}}_{D} &=  \vec{g} - \eta_v \vec{v}_{D} + \frac{F}{m} \mathbf{A} \cdot \vec{e}_{z} = f_v(x_D) +\frac{1}{m} \mathbf{A} \cdot \vec{e}_{z} \cdot F, \\
    \dot{\vec{\omega}}^B &= -\mathbf{J}^{-1} \cdot (\vec{\omega}^B \times \mathbf{J} \cdot \vec{\omega}^B) - \eta_{\omega} + \mathbf{J}^{-1} \cdot \vec{\tau} = f_{\omega}(x_D)+\mathbf{J}^{-1} \cdot \vec{\tau}, \\
    \text{vec}(\dot{\mathbf{A}}) &= \text{vec}(-\vec{\omega}^B \times \mathbf{A}) = f_A(x_D). \\
\end{align*}

The drone's differential dynamics can be defined by the control-affine form $\dot{x}_{D} = f_D(x_D)+ \mathbf{C} \cdot u_D$, where

\begin{equation*}
    f_R(x_D) = \begin{bmatrix}
        f_p(x_D) \\
        f_v(x_D) \\
        f_{\omega}(x_D) \\
        f_A(x_D) \\
    \end{bmatrix}, \quad 
    \mathbf{C} = \begin{bmatrix}
        \mathbf{0} & \mathbf{0} \\
        \frac{1}{m} \mathbf{A} \cdot \vec{e}_{z} & \mathbf{0} \\
        \mathbf{0} & \mathbf{J}^{-1} \\
        \mathbf{0} & \mathbf{0} \\
    \end{bmatrix}.
\end{equation*}

\begin{comment}
The drone's dynamics are simplified as a double integrator in three dimensions. They're given by:

\begin{align*}
    \dot{x}_{D} = f_D(x_D, u_D) = \begin{bmatrix}
        f_p(x_D, u_D) \\
        f_v(x_D, u_D) \\
    \end{bmatrix} = 
    \begin{bmatrix}
        \vec{v_D} \\
        \vec{a_D} + \vec{g} 
    \end{bmatrix} =
    \begin{bmatrix}
        \dot{\vec{p_D}} \\
        \dot{\vec{v_D}} \\
    \end{bmatrix}
\end{align*}

In which:
\begin{description}
    \item[$\vec{p}_{D} \in \mathbb{R}^3$] is the position of the center of mass, expressed in $\mathcal{F}_I$.
    \item[$\vec{v}_{D} \in \mathbb{R}^3$] is the linear speed of the center of mass, expressed in $\mathcal{F}_I$. \\
\end{description}

Additional variables are defined as:
\begin{description}
    \item[$\vec{a}_D \in \mathbb{R}^3$] is the thruster acceleration vector acting on the rocket, expressed in $\mathcal{F}_B$.
    \item[$\vec{g} \in \mathbb{R}^3$] is the gravitational acceleration vector, expressed in $\mathcal{F}_I$. \\
\end{description}

The input vector $u_D \in \mathbb{R}^3$ is defined as
\begin{equation*}
    u_D = \vec{a}_D
\end{equation*}
\end{comment}

\section{Missile Defense}

The scenario in study  can be thought of a defense-intercept game with two asymmetric players, with each player having their hierarchy of objectives to satisfy.

The missile's hierarchy of objectives are static target interception and dynamic pursuer evasion. The drone's hierarchy of objectives are dynamic target interception and static area defense. \\

Each objective poses its own control problem, but each player's GNC structure must account for all its objectives. Regarding a player's navigation system, it must be able to estimate its own and the opponent's state relative to each other and relative to the critical area's position. As for the guidance system, it must apply decision-making based on the player's navigation data and it's goals to plan the most optimal trajectory. Finally, the flight control system calculates the appropriate actuation commands given the player's dynamics, based on its navigation and guidance data. \\

Figure~\ref{fig:intercept_diag} illustrates both the missile's and drone's expected behaviors in a successful defense scenario. \textit{C} identifies the critical structure, \textit{D} identifies the drone, \textit{R} identifies the missile, $\vec{r}_R$ is the missile's best feasible planned trajectory, $\vec{d}_{DR}$ is the distance vector between the missile and the drone, and $\vec{d}_{RC}$ is the distance vector between the missile and the critical structure. From instant \textit{k} to instant \textit{k+1} the distance vectors shrink, and the missile's planned trajectory can no longer hit the target thanks to its collision avoidance.

\begin{figure}[htbp]
  \centering
    \subbottom[At instant k]{%
    \label{fig:intercept_diag0}
    \includegraphics[width=0.45\textwidth]{ch2_intercept_diag0}
    }
    \subbottom[At instant k+1]{%
    \label{fig:intercept_diag1}
    \includegraphics[width=0.45\textwidth]{ch2_intercept_diag1}
    }
  \caption{Missile defense evolution in time}
  \label{fig:intercept_diag}
\end{figure}

For the flight control system, it is necessary to define a control law. Since either player is described by a non-linear model, several methods may be used, such as model linearization around an ideal trajectory or NMPC. Optimal Control-based methods, which obtain the optimal input actuation by minimizing some cost function related to the objective, stand out for their ability to join guidance with control, as they do not need explicit path planning and trajectory generation to generate the optimal input. These characteristics favor the control problems at hand, seeing that players do not need to follow specific trajectories, only to converge on their targets while satisfying their other objectives. CLF/CBF-type controllers are also well suited to these problems, since the evasion and defense objectives are essentialy barrier conditions. \\

For the scenario in which the rocket does not evade the drone, let $V_R : \mathcal{X} \rightarrow \mathbb{R}$ a quadratic function representing the distance between the rocket and the critical structure. The rocket's objective can be represented by the optimization problem:
\begin{align*}
    \underset{u}{\text{min}} &\quad V_{R}(x,u) \\
    \text{s.t.:} &\quad \dot{x} = f_R(x, u) \\
    &\quad x \in \mathcal{X}_{R} \\
    &\quad u \in \mathcal{U}_R \\
\end{align*}

The drone's objective is dual; to intercept the rocket as well as diverting the rocket away from the critical structure. Let $V_D : \mathcal{X} \rightarrow \mathbb{R}$ a quadratic function representing the distance between the drone and the rocket, and $V_{\text{crit}} : \mathcal{Z} \rightarrow \mathbb{R}$ a quadratic CBF representing the intersection between the reachable set of the rocket and the critical area. Furthermore, consider $\gamma (\cdot)$ is a $\mathcal{K}$ class function. The drone's objective can be represented by the optimization problem:
\begin{align*}
    \underset{u}{\text{min}} &\quad V_{D}(x,z,u) \\
    \text{s.t.:} &\quad \dot{x} = f_D(x, u) \\
    &\quad \dot{V}_{\text{crit}}(z) \le -\gamma (V_{crit}(z)) \\
    &\quad x \in \mathcal{X}_{D} \\
    &\quad z \in \mathcal{X}_R \\
    &\quad u \in \mathcal{U}_D \\
\end{align*}

